\documentclass[aspectratio=169,11pt]{beamer}
\usepackage{teaching_slides}



\title [MPEC]{MPEC Approach}
\author{C.Conlon - thanks to late great Che-Lin Su}
\institute{Grad IO }
\date{}
\setbeamerfont{equation}{size=\tiny}

\begin{document}
\begin{frame}
\titlepage
\end{frame}



\frame{\frametitle{Extremum Estimators}
Often faced with extremum estimator problems in econometrics (ML, GMM, MD, etc.)  that look like:
\begin{align*}
\hat{\theta} = \arg \max_{\theta} Q_n(\theta), \quad  \theta \in \Theta
\end{align*}
Many economic problems contain constraints, such as: market clearing (supply equals demand), consumers consume their entire budget set, or firms' first order conditions are satisfied.  A natural way to represent these problems is as constrained optimization.
}

\begin{frame}{Constrained Problems}
\begin{block}{MPEC}
\begin{align*}
\hat{\theta} = \arg \max_{\alert{\theta},\textcolor{blue}{P}} Q_n(\alert{\theta},\textcolor{blue}{P}), \quad  \mbox{ s .t. } \quad \Psi(\textcolor{blue}{P},\alert{\theta}) =0,  \quad \alert{\theta} \in \Theta
\end{align*} 
\end{block}
\pause
\begin{block}{Fixed Point / Implicit Solution}
In much of the literature the tradition has been to express the solutions $\Psi(P,\theta) =0$ implicitly as $P(\theta)$:
\begin{align*}
\hat{\theta} = \arg \max_{\theta} Q_n(\alert{\theta},\textcolor{blue}{P}(\alert{\theta})),  \quad \alert{\theta} \in \Theta
\end{align*}
\end{block}
\end{frame}

\begin{frame}{NFXP vs MPEC}
Probably you were taught some things that weren't quite right
\begin{itemize}
\item Fewer parameters $\rightarrow$ easier problems to solve!
\item Reformulate problems with fixed points or implicit solutions to \alert{concentrate out} parameters.
\item But sometimes this makes the problem more complicated (saddle points, complicated Hessians, etc.)
\end{itemize}
MPEC says do the opposite:
\begin{itemize}
\item Add lots of parameters
\item But add them with simple constraints (linear or quadratic).
\item Idea: Make the Hessian as close to constant, block diagonal, sparse, etc. as possible.
\end{itemize}
Mostly this is in response to change in technology: nonlinear solvers supporting large sparse Hessians.
\end{frame}



\begin{frame}{Rust Problem}
\footnotesize
\begin{itemize}
\item Bus repairman sees mileage $x_t$ at time $t$ since last overhaul
\item Repairman chooses between overhaul and normal maintenance
\begin{align*}
u(x_t, d_t; \alert{\theta^c}, \alert{RC}) =
\begin{cases} -c(x_t;\alert{\theta^c}) \quad &\mbox{if} \quad d_t=0\\
-(\alert{RC} + c(0;\alert{\theta^c})) &\mbox{if} \quad d_t =1\\
\end{cases}
\end{align*}
\item Repairman solves DP:
\begin{align*}
V_{\alert{\theta}} (x_t, \varepsilon_t) = \max_{\{d_t,d_{t+1},\ldots\}} E \left\{  \sum_{j=t}^{\infty} \beta^{j-t}   [u(x_j,d_j;\alert{\theta}) + \varepsilon_{j}(d_j)]| x_t, \varepsilon_t \right\}
\end{align*}
\end{itemize}
\begin{itemize}
%\item Observe mileage $x_t$ and decision $d_t$ but not cost.
%\item Assumes extreme value distribution for $\varepsilon_t(d_t)$
\item Structural parameters to be estimated $\alert{\theta = (\theta^c, RC, \theta^f)}$.
\item Coefficients of cost function $c(x;\alert{\theta^c}) = \theta^c_1 x + \theta^c_2 x^2$
\item Transition probabilities in mileages $f(x_{t+1} | x_t, d_t; \alert{\theta^f})$
\end{itemize}
\end{frame}


\begin{frame}{Rust Problem}
\begin{itemize}
\item Data: time series $(x_t,d_t)_{t=1}^T$
\item Likelihood function
\begin{align*}
\mathcal{L}(\theta) &= \prod_{t=2}^T P(d_t | x_t; \alert{\theta^c}, \alert{RC}) f(x_t | x_{t-1}, d_{t-1};\alert{\theta^f})\\
\mbox{ with }  P(d | x; \alert{\theta^c}, \alert{RC}) &= \frac{\exp[u(x,d;\alert{\theta^c},\alert{RC}) + \beta \textcolor{blue}{EV}_{\alert{\theta}}(x,d)]}
{\sum_{d' \in \{0,1\}} \exp[u(x,d';\alert{\theta^c},\alert{RC}) + \beta \textcolor{blue}{EV}_{\alert{\theta}}(x,d')]}\\
\textcolor{blue}{EV}_{\alert{\theta}}(x,d) &= T_{\alert{\theta}}(\textcolor{blue}{EV}_{\alert{\theta}})(x,d)\\
\end{align*}
\vspace{-1cm}
\begin{align*}
\equiv \int_{x'=0}^{\infty} &\log \left[ \sum_{d' \in \{0,1\}} \exp[u(x',d';\alert{\theta^c},\alert{RC}) + \beta \textcolor{blue}{EV}_{\alert{\theta}}(x',d')]  \right] f(dx' | x,d;\alert{\theta^f})
\end{align*}
\end{itemize}
\end{frame}


\begin{frame}{Rust Problem}
\begin{itemize}
\item Outer Loop: Solve Likelihood
\begin{align*}
\max_{\alert{\theta} \geq 0} \mathcal{L}(\theta) &= \prod_{t=2}^T P(d_t | x_t; \alert{\theta^c}, \alert{RC}) f(x_t | x_{t-1}, d_{t-1};\alert{\theta^f})\\
\end{align*}
\item Convergence test: $\norm{\nabla_{\theta} \mathcal{L}(\alert{\theta})} \leq \alert{\epsilon_{out}}$
\item Inner Loop: Compute expected value function $ \textcolor{blue}{EV}_{\alert{\theta}}$ for a given $\alert{\theta}$
\item $ \textcolor{blue}{EV}_{\alert{\theta}}$ is the implicit expected value function defined by the Bellman equation or the fixed point function
\begin{align*}
 \textcolor{blue}{EV}_{\alert{\theta}} = T_{\alert{\theta}}( \textcolor{blue}{EV}_{\alert{\theta}})
\end{align*}
\item Convergence test: $\norm{ EV^{(k+1)}_{\alert{\theta}} - EV^{(k)}_{\alert{\theta}} } \leq \alert{\epsilon_{in}}$
\item Start with contraction iterations and polish with Newton Steps
\end{itemize}
\end{frame}

\begin{frame}{NFXP Concerns}
\begin{itemize}
\item Inner-loop error propagates into outer-loop function and derivatives
\item NFXP needs to solve inner-loop exactly each stage of parameter search
\begin{itemize}
\item to accurately compute the search direction for the outer loop
\item to accurately evaluate derivatives for the outer loop
\item for outer loop to converge!
\end{itemize}
\item Stopping rules: choosing inner-loop and outer-loop tolerance
\begin{itemize}
\item inner loop can be slow: contraction mapping is linearly convergent
\item tempting to loosen inner loop tolerance $\alert{\epsilon_{in}}$ (such as $1e-6$ or larger!).
\item Outer loop may not converge with loose inner loop tolerance.
\begin{itemize}
\item check solver output message
\item tempting to loosen outer loop tolerance $\alert{\epsilon_{out}}$ to promote convergence ($1\text{e-}3$ or larger!).
\end{itemize}
\end{itemize}
\end{itemize}
\end{frame}

\begin{frame}{Convergence Properties (Su and Judd)}
\begin{itemize}
\item $\mathcal{L}(\textcolor{blue}{EV}(\alert{\theta},\alert{\epsilon_{in}} ), \alert{\theta} )$ the programmed outer loop objective function
\item $L$: the Lipschitz constant (like modulus) of inner-loop contraction mapping
\item Analytic derivatives $\nabla_{\theta} \mathcal{L}(\textcolor{blue}{EV}(\alert{\theta},\alert{\epsilon_{in}} ), \alert{\theta} )$ is provided: $\alert{\epsilon_{out}} = O(\frac{L}{1-L} \alert{\epsilon_{in}})$
\item Finite-difference derivatives are used: $\alert{\epsilon_{out}} = O ( \sqrt{\frac{L}{1-L}} \alert{\epsilon_{in}} )$ 
\end{itemize}
\end{frame}

\begin{frame}{MPEC Alternative (Su and Judd)}
\begin{itemize}
\item Form the augmented likelihood function for data $X = (x_t, d_t)_{t=1}^T$
\begin{align*}
\mathcal{L}(\textcolor{blue}{EV}, \alert{\theta}; X) &= \prod_{t=2}^T P(d_t | x_t; \alert{\theta^c}, \alert{RC}) f(x_t | x_{t-1}, d_{t-1};\alert{\theta^f})\\
\mbox{ with }  P(d | x; \alert{\theta^c}, \alert{RC}) &= \frac{\exp[u(x,d;\alert{\theta^c},\alert{RC}) + \beta \textcolor{blue}{EV}(x,d)]}
{\sum_{d' \in \{0,1\}} \exp[u(x,d';\alert{\theta^c},\alert{RC}) + \beta \textcolor{blue}{EV}(x,d')]}\\
\end{align*}
\item Rationality and Bellman equation imposes a relationship between $\alert{\theta}$ and $\textcolor{blue}{EV}$
\begin{align*}
\textcolor{blue}{EV} = T(\textcolor{blue}{EV},\alert{\theta})
\end{align*}
\item Solve constrained optimization problem
\begin{align*}
\max_{(\alert{\theta},\textcolor{blue}{EV})} \mathcal{L}(\textcolor{blue}{EV}, \alert{\theta}; X)\\
\mbox{ subject to }  \textcolor{blue}{EV} = T(\textcolor{blue}{EV},\alert{\theta})
\end{align*}
\end{itemize}
\end{frame}


\begin{frame}{MPEC Alternative$^2$}
The previous approach solves the problem ``on the grid''.
\begin{itemize}
\item $x_t$ takes on discrete values.
\item We only evaluate $EV(x_t,d_t)$ at values on the grid.
\item We don't evaluate $EV(x_t,d_t)$ and values between $[x_s,x_{s+1}]$.
\end{itemize}
If we did we would have to \alert{interpolate}.
\begin{itemize}
\item Macroeconomists tend to use \alert{cubic splines}
\item Could use global orthogonal polynomials $EV(x) \approx a_0 + a_1 b_1(x)+ a_2 b_2(x) + \ldots$
\begin{itemize}
\item Advantage is that solving polynomial equation for $\textcolor{blue}{EV} = T(\textcolor{blue}{EV},\alert{\theta})$ is pretty easy.
\item Can easily switch to continuous state space without any additional complications:  $f(x_{t+1} | x_{t})$ must also be continuous.
\end{itemize}
\end{itemize}
\end{frame}

\begin{frame}{Monte Carlo: Su and Judd's design, our code}
\footnotesize
\begin{itemize}
\item Rust's bus model as in Su and Judd (2012): $N = 175$ mileage states, $c(x;\theta^c) = 0.001\,\theta^c_1\,x$ (their $\theta_{11}$), $u(x,1) = -RC$; truth $(\theta^c_1, RC) = (2.4569, 11.7257)$, $\theta^f = (.0937, .4475, .4459, .0127, .0002)$.
\item For each $\beta \in \{0.975, 0.985, 0.995, 0.999, 0.9995, 0.9999\}$: 250 datasets of 50 buses $\times$ 120 months (6{,}000 obs.), 5 random starts each ($\theta^c_1 \sim U[0.5,5]$, $RC \sim U[5,20]$): 1{,}250 runs per method.
\item Two-step: $\hat\theta^f$ from frequencies, then $(\theta^c_1, RC)$ by three estimators of the same partial likelihood:
\begin{itemize}
\item \alert{NFXP-SA}: the inner loop is successive approximations until $\|EV_{k+1} - EV_k\| < 10^{-10}$ (Su and Judd's NFXP).
\item \alert{NFXP-NK}: Rust's polyalgorithm: successive approximations until the contraction rate is close to $\beta$, then Newton--Kantorovich. Outer loop: trust region on the BHHH matrix; gradient by the implicit function theorem.
\item \alert{MPEC}: $\max_{\theta, EV} \ell(\theta, EV)$ s.t.\ $EV = T_\theta(EV)$, SLSQP with exact Jacobians, $EV$ rescaled by $(1-\beta)$.
\end{itemize}
\item Python (NumPy/SciPy; the successive-approximation loop compiled with JAX), one core per run. \alert{Converged} = the solver reports success \emph{and} reaches the best likelihood found on that dataset.
\end{itemize}
\end{frame}

\begin{frame}{Monte Carlo results: speed}
\scriptsize
\begin{center}
\begin{tabular}{l rrr c rrr c rrr}
\toprule
 & \multicolumn{3}{c}{MPEC} & & \multicolumn{3}{c}{NFXP-NK} & & \multicolumn{3}{c}{NFXP-SA}\\
\cmidrule{2-4} \cmidrule{6-8} \cmidrule{10-12}
$\beta$ & Conv. & CPU (ms) & Iter. & & Conv. & CPU (ms) & SA + NK & & Conv. & CPU (ms) & SA steps\\
\midrule
0.975  & 1{,}250 & 36.3 & 26.9 & & 1{,}250 & 10.1 & 93 + 31  & & 1{,}250 & 50.1    & 6{,}484\\
0.985  & 1{,}250 & 36.6 & 27.2 & & 1{,}250 & 10.3 & 97 + 30  & & 1{,}250 & 78.9    & 10{,}293\\
0.995  & 1{,}250 & 34.4 & 27.9 & & 1{,}250 &  9.5 & 100 + 29 & & 1{,}250 & 215.6   & 29{,}708\\
0.999  & 1{,}250 & 41.3 & 28.3 & & 1{,}250 &  9.1 & 97 + 28  & & 1{,}250 & 926.5   & 143{,}205\\
0.9995 & 1{,}250 & 38.9 & 28.5 & & 1{,}250 &  8.9 & 95 + 27  & & 1{,}250 & 1{,}815.2 & 283{,}658\\
0.9999 & 1{,}249 & 30.5 & 28.4 & & 1{,}250 &  7.9 & 100 + 27 & & 1{,}250 & 9{,}100.6 & 1{,}447{,}801\\
\bottomrule
\end{tabular}
\end{center}
Conv.: converged runs out of 1{,}250. CPU: mean milliseconds per estimation over converged runs, one core. Iter.: SLSQP iterations. SA + NK / SA steps: inner-loop iterations per estimation, summed over the outer loop.
\end{frame}

\begin{frame}{Monte Carlo results: estimates}
\footnotesize
\begin{center}
\begin{tabular}{l cc}
\toprule
 & $\theta^c_1$ & $RC$\\
\midrule
truth & 2.4569 & 11.7257\\
\midrule
$\beta = 0.975$  & 2.597 (0.474) & 12.225 (1.537)\\
$\beta = 0.985$  & 2.587 (0.453) & 12.173 (1.471)\\
$\beta = 0.995$  & 2.573 (0.478) & 12.058 (1.531)\\
$\beta = 0.999$  & 2.551 (0.423) & 12.022 (1.318)\\
$\beta = 0.9995$ & 2.595 (0.471) & 12.172 (1.477)\\
$\beta = 0.9999$ & 2.549 (0.430) & 12.010 (1.276)\\
\bottomrule
\end{tabular}
\end{center}
Mean (standard deviation) across the 250 datasets of the estimate at the best likelihood. Conditional on convergence the three methods return the same estimates up to the solvers' stopping tolerances: on any dataset they differ by at most 0.011 in $\theta^c_1$ and 0.04 in $RC$, no more than different starts of one method and under 3\% of a standard deviation. The upward bias in both parameters (a quarter of a standard deviation) is finite-sample bias of the estimator, not of any one algorithm.
\end{frame}

\begin{frame}{What changed between 2012 and 2016?}
\begin{itemize}
\item Su and Judd (2012): ``MPEC is significantly faster than NFXP.'' But their NFXP solved the inner loop by successive approximations, which takes about $\ln(\text{tol}/\|EV_0 - EV\|)/\ln\beta$ steps per solve: thousands per estimation at $\beta = 0.975$, over a million at $\beta = 0.9999$.
\item Iskhakov, Lee, Rust, Schjerning and Seo (Ecma 2016): with Rust's own NFXP-NK polyalgorithm, NFXP is about as fast as MPEC. In our runs each inner solve takes about a dozen successive approximations and 3--4 Newton steps at every $\beta$, and NFXP-NK is about four times faster than our SLSQP-based MPEC.
\item Lessons:
\begin{itemize}
\item Compare good implementations: the inner solver matters more than nested vs.\ constrained.
\item Exact derivatives do the heavy lifting: implicit differentiation for NFXP, exact Jacobians for MPEC.
\item MPEC is sensitive to scaling: $EV = O(1/(1-\beta))$. Without rescaling $EV$, MPEC converges in every run for $\beta \le 0.995$ (in twice as many iterations), in 76 of 1{,}250 runs at $\beta = 0.999$, and in none for $\beta \ge 0.9995$.
\end{itemize}
\end{itemize}
\end{frame}

\frame{\frametitle{BLP Demand Example}
\begin{exampleblock}{BLP 1995}
The estimator solves the following mathematical program:
\begin{align*}
\min_{\alert{\theta_2}}\, & g(\xi(\alert{\theta_2}))' W g(\xi(\alert{\theta_2})) \quad \mbox{ s.t. } \\
g(\xi(\alert{\theta_2})) &= \frac{1}{N} \sum_{\forall j,t} \xi_{jt} (\alert{\theta_2}) ' z_{jt} \\
\xi_{jt}(\alert{\theta_2}) &= \delta_{jt}(\alert{\theta_2}) - x_{jt} \beta + \alpha p_{jt} \\
\sigma_{jt}(\delta_t(\alert{\theta_2}),\alert{\theta_2}) &= \int \frac{\exp[\delta_{jt}(\alert{\theta_2})+ \mu_{ijt}]}{1+\sum_k \exp[\delta_{kt}(\alert{\theta_2})+ \mu_{ikt}]} f(\mu | \alert{\theta_2})\\
\log s_{jt}  &= \log \sigma_{jt}(\delta_t(\alert{\theta_2}),\alert{\theta_2}) \quad \forall j,t
\end{align*}
\end{exampleblock}
}

\frame{\frametitle{BLP Algorithm}
The estimation algorithm is generally as follows:
\begin{enumerate}
\item Guess a value of nonlinear parameters $\alert{\theta_2}$
\item Compute $\sigma_{jt}(\delta_t,\alert{\theta_2})$ via integration
\item Iterate on $\delta_{jt}^{h+1}  = \delta_{jt}^h +  \log s_{jt} - \log \sigma_{jt}(\delta^{h}_t,\alert{\theta_2})$ to find the $\delta$ that satisfies the share equation
\item IV Regression $\delta$ on observable $X$ and instruments $Z$ to get residual $\xi$.
\item Use $\xi$ to construct $g(\xi(\alert{\theta_2}))$.
\item Possibly construct other errors/instruments from supply side.
\item Construct GMM Objective
\end{enumerate}
The idea is that $\delta(\alert{\theta_2})$ is an implicit function of the nonlinear parameters $\alert{\theta_2}$. And for each guess we find that implicit solution for reduce the parameter space of the problem. But the Jacobian: $\frac{\partial \symbf{\xi_t}}{\partial \theta_2}(\theta_2) = -{\left[\frac{\partial \symbf{\sigma}_t}{\partial \symbf{\delta_t}}(\theta_2)\right]^{-1} \left[\frac{\partial \symbf{\sigma}_t}{\partial \theta_2}(\theta_2)\right]} $ is complicated to compute.
}



\frame{\frametitle{Dube Fox Su 2012}
\footnotesize
\begin{exampleblock}{BLP-MPEC}
The estimator solves the following mathematical program:
\begin{align*}
\label{blpmpec}
\min_{\alert{\Sigma,\alpha,\beta, \xi}}  & g(\alert{\xi})' W g(\alert{\xi}) \quad \mbox{ s.t. } \\
g(\alert{\xi}) &= \frac{1}{N} \sum_{\forall j,t} \alert{\xi_{jt}'} z_{jt} \\
\sigma_{jt}(\alert{\Sigma,\alpha,\beta, \xi}) &=  \sum_i w_i \frac{\exp[x_{jt} \alert{\beta} + \alert{\xi_{jt}} - \alert{\alpha} p_{jt} +  x_{jt}\alert{\Sigma} \nu_{i} ] }
{1+ \sum_k \exp[x_{kt} \alert{\beta} + \alert{\xi_{kt}} - \alert{\alpha} p_{kt} + x_{kt}\alert{\Sigma} \nu_{i} ]} \\
\log s_{jt}  &= \log \sigma_{jt}(\alert{\Sigma,\alpha,\beta, \xi}) \quad \forall j,t
\end{align*}
\end{exampleblock}
\begin{itemize}
\item Expand the parameter space of the nonlinear search to include $\alert{\alpha,\beta,\xi}$
\item Don't have to solve for $\xi$ except at the end.
\item No implicit functions of $\theta_2$ and $\left[\frac{\partial \symbf{\sigma}_t}{\partial \Sigma}(\alert{\Sigma,\alpha,\beta, \xi})\right]$ is straightforward (no matrix inverse!).
\item Sparsity!
\end{itemize}
}




\frame{\frametitle{Nevo Results}
\begin{center}
\rowcolors[]{1}{RoyalBlue!5}{RoyalBlue!15} 
\begin{tabular}{lrrrr}
& Nevo & BLP-MPEC & EL\\
Price& -28.189  & -62.726  & -61.433  \\
$\Sigma_{p}$ &0.330 & 0.558 & 0.524\\
$\Sigma_{const}$ &2.453 &3.313 & 3.143\\
$\Sigma_{sugar} $& 0.016&-0.006 & 0\\
$\Sigma_{mushy} $&0.244 &0.093&0.085\\
$\Pi_{p,inc}$ &15.894  &588.206 & 564.262\\
$\Pi_{p,inc2}$ &-1.200  &-30.185 &-28.930 \\
$\Pi_{p,kid}$ &2.634 &11.058 &   11.700\\
$\Pi_{c,inc}$ &5.482 & 2.29084&2.246\\
$\Pi_{c,age} $& 0.2037&1.284 & 1.37873 \\
GMM &29.3611 &4.564 & \\
EL & & &-17422  \\
Time & 28 s & 12s & 19s  \\
\end{tabular}
\end{center}
}

\begin{frame}{Another Example: Empirical Likelihood}
Consider the GMM moment condition $\E[g(Z_i,\theta)]=0$
\begin{itemize}
\item A common example is $\E[z_i' (y_i-x_i \beta)]=0$
\item Now consider re-weighting the data with $\rho_i$ so that $\sum_i \rho_i g(Z_i,\theta)=0$
\begin{itemize}
\item We also need that $\sum_i \rho_i =1$ and $\rho_i \in (0,1)$ (ie: $\rho$ is a valid probability measure)
\item Idea: penalize distance between $\rho_i$ and $\frac{1}{n}$ with $f(\rho)$
\end{itemize}
\item Idea: your moments always hold (because economic theory)
\item How much would you have to re-sample data?
\end{itemize}
\end{frame}



\begin{frame}{Another Example: (Generalized) Empirical Likelihood}
\begin{align*}
\min_{\alert{\theta}, \textcolor{blue}{\rho} }\sum_i f(\textcolor{blue}{\rho_i}) \quad &\text{ s.t } \sum_i \textcolor{blue}{\rho_i} \cdot g(Z_i,\alert{\theta})= 0\\
& \sum_i \textcolor{blue}{\rho_i} =1 \quad \textcolor{blue}{\rho_i} \geq 0
\end{align*}
Choices of $f(\rho_i)$:
\begin{itemize}
\item Empirical Likelihood: $f(\rho_i) = - \log \rho_i$ [This is NPMLE]
\item Continuously Updating GMM (CUE): $f(\rho_i) = \rho_i^2$ [Derivation is \alert{not} obvious]
\item Exponential Tilt: $f(\rho_i) = \rho_i \log \rho_i $ [has some ``robustness'' properties]
\item Adversarial Method of Moments: Conlon, Kim, Manresa, Li (2025)
\item See Kitamura's Handbook Chapter for more details (and convex analysis) or Newey Smith (2004) for asymptotic bias.
\end{itemize}
\end{frame}


\begin{frame}{Empirical Likelihood: NFXP Solution}
Start with the Lagrangian:
\begin{align*}
\mathcal{L} &=\sum_{i=1}^{n} \log \rho_{i}+\lambda\left(1-\sum_{i=1}^{n} \rho_{i}\right)-n \gamma^{\prime} \sum_{i=1}^{n} \rho_{i}\cdot g\left(Z_{i}, \theta\right)
\end{align*}
Take derivatives with respect to \textcolor{blue}{$\rho$} and $\alert{\theta}$ and Lagrange multipliers to get $\textcolor{blue}{\gamma}$ and concentrate out $\textcolor{blue}{\rho_i}(\alert{\theta})$
\begin{align*}
\textcolor{blue}{\hat{\gamma}}(\alert{\theta}) &=\arg \min\, _{\textcolor{blue}{\gamma} \in \mathbb{R}^{q}}-
\sum_{i=1}^{n} \log \left(1+\textcolor{blue}{\gamma}^{\prime} g\left(Z_{i}, \alert{\theta}\right)\right) \\
\textcolor{blue}{\hat{\rho}}_{i}(\alert{\theta}) &=\frac{1}{n\left(1+\textcolor{blue}{\hat{\gamma}}(\alert{\theta})^{\prime} g\left(Z_{i}, \alert{\theta}\right)\right)}, \quad \hat{\lambda}=n
\end{align*}
The resulting dual is a saddle-point problem:
\begin{align*}
\hat{\theta}_{E L}=\arg \max _{\alert{\theta} \in \Theta} \ell_{N P}(\alert{\theta})=\arg \max _{\alert{\theta} \in \Theta} \min\, _{\textcolor{blue}{\gamma} \in \mathbb{R}^{q}}-\sum_{i=1}^{n} \log \left(1+\textcolor{blue}{\gamma}^{\prime} g\left(Z_{i}, \alert{\theta}\right)\right)
\end{align*}
But the primal (MPEC) problem is much easier... unless $g(Z_i,\theta)$ is trivial.
\end{frame}

\section{Thanks!}

\end{document}


